% Source: https://artofproblemsolving.com/community/c3h1328182_geometry_length_chasing
% Stable identity: aops-topic-1328182-q01
Place the triangle with vertices $P=(0,0)$, $Q=(6,0)$, $R=(3,3\sqrt3)$. Each square has two outer corners, obtained by moving the ends of its side $6$ units along the outward normal of that side. Comparing the distances between all nine vertices, the farthest pair is formed by the outer corner of the square on $RP$ at $P$ and the outer corner of the square on $QR$ at $Q$. The outward unit normals of $RP$ and $QR$ are $\left(-\dfrac{\sqrt3}{2},\dfrac{1}{2}\right)$ and $\left(\dfrac{\sqrt3}{2},\dfrac{1}{2}\right)$, so these corners are
\[
A=\left(-3\sqrt3,\ 3\right),\qquad B=\left(6+3\sqrt3,\ 3\right).
\]
Their distance is $6+6\sqrt3=6+\sqrt{108}$.

\begin{asy}
usepackage("amsmath");
size(5cm);
defaultpen(fontsize(10pt));
pair A=(0,0),B=(6,0),C=(3,3*sqrt(3));
draw(A--B--C--cycle,linewidth(0.9));
pair[] P={A,B,C}; for(int i=0;i<3;++i){pair U=P[i],V=P[(i+1)%3]; pair n=rotate(-90)*(V-U); draw(U--U+n--V+n--V,linewidth(0.9));}
pair X=A+rotate(-90)*(A-C), Y=B+rotate(-90)*(C-B);
draw(X--Y,red+dashed); label("$A$",X,W); label("$B$",Y,E);
dot(X);dot(Y);
label("\mbox{}",(min(currentpicture,user=true)+max(currentpicture,user=true))/2,black); // keeps the text after the figure black when LaTeX typesets the labels inline
\end{asy}

So $m=6$, $n=108$, and $m+n=\boxed{114}$.
